% Осень26!B19; нумерация внутри столбца без пустых строк.
% Источник: лекция 3, §4 (строка 115)
\definitionticket{16}{Сигма-конечная мера}
\label{definition:16}

\begin{definition}
Мера $\mu$ на $\sigma$-алгебре $\mathcal A$ подмножеств $X$ называется \textbf{$\sigma$-конечной}, если существуют множества $A_n\in\mathcal A$ такие, что
\[
  X=\bigcup_{n=1}^{\infty}A_n,
  \qquad \mu(A_n)<+\infty\quad(n\in\mathbb N).
\]
\end{definition}
